\subsection{Kirchhoffsche Gesetze 4}
\Aufgabe{
    Gegeben ist die dargestellte Schaltung aus den drei Widerständen $R_1, R_2$ und $R_3$ und
    den beiden Quellen $U_{\mathrm{Q}}$ und $I_{\mathrm{Q}}$. \newline
    
    \begin{center}
    \input{Tikz/src/Aufgaben/Kirchhoff_A4}
    \end{center}
    \begin {enumerate}[label=\alph*)]
        \item Bestimmen Sie den Innenwiderstand der Schaltung bezüglich der Klemmen A und B.
        \item Berechnen Sie die Leerlaufspannung $U_{\mathrm{ABL}}$ zwischen den Klemmen A und B.
        \item Geben Sie die Ersatzspannungsquelle für die Schaltung nach dem vorstehenden Bild an und kennzeichne die Größen.
        \item Welche Verlustleistung $P_{\mathrm{V}}$ wird in der Schaltung nach dem vorstehenden Bild umgesetzt?
        \item Die Klemmen A und B werden kurzgeschlossen. Welcher Strom $I_{\mathrm{K}}$ fließt durch diese Verbindung
        zwischen den Klemmen A und B?
        \item Welche Spannung $U_{\mathrm{AB}}$ stellt sich ein, wenn an die Klemmen A und B ein Widerstand $R$ angeschlossen wird? $(R_1 = R_2 = R_3 = R)$    
    \end{enumerate}
}

\Loesung{



    \begin {enumerate}[label=\alph*)]
        \item Bestimmen Sie den Innenwiderstand der Schaltung bezüglich der Klemmen A und B.
        

        Zum bestimmen des Innenwiderstandes: 
         \begin{itemize}
            \item Stromquellen öffnen
            \item Spannungsquellen kurzschließen
        \end{itemize}

        \begin{center}
            \begin{tikzpicture}
    \draw(3,3) to[short,o-](2,3)
    to[R,l_=$R_1$](0,3)
    to[short](0,0)
    to[short](5,0)
    to[R,l=$R_3$](5,3)
    to[short,-o](4,3);
    \draw(2,0)to[R,l=$R_2$](2,3);
    
    \node at (3.2,3.3){A};
    \node at (3.8,3.3){B};
\end{tikzpicture}
        \end{center}

    
    $ R_\mathrm{i,AB}= R_3 + \frac{R_1 \cdot R_2}{R_1 + R_2}$
        


        \item Berechnen Sie die Leerlaufspannung $U_{\mathrm{ABL}}$ zwischen den Klemmen A und B.
        
        Zur Bestimmung der Leerlaufspannung oft sinnvoll: Stromquelle in Spannungsquelle umrechnen:
        \begin{itemize}
            \item Parallelgeschalteten Innenwiderstand in Reihe schalten
            \item Spannungspfeil der Ersatzspannungsquelle entgegen dem ursprünglichen Strompfeil zeichnen
        \end{itemize}

        \begin{center}
            \input{Tikz/src/Aufgaben/Kirchhoff_A4_lsg2}
        \end{center}

            $I_\mathrm{R3} = 0 \rightarrow U_\mathrm{R3} = 0$\

            $U_\mathrm{AB,L} = U_\mathrm{Q} \cdot \frac{R_2}{R_1+R_2} + I_\mathrm{Q} \cdot R_3$



        
        \item Geben Sie die Ersatzspannungsquelle für die Schaltung nach dem vorstehenden Bild an und kennzeichne die Größen.
        
        \begin{center}
            \input{Tikz/src/Aufgaben/Kirchhoff_A4_lsg3}
        \end{center}

        $U_0 = U_\mathrm{AB,L} =  U_\mathrm{Q} \cdot \frac{R_2}{R_1+R_2} + I_\mathrm{Q} \cdot R_3$\\

        $ R_\mathrm{i} = R_\mathrm{i,AB}= R_3 + \frac{R_1 \cdot R_2}{R_1 + R_2}$

        

        \item Welche Verlustleistung $P_{\mathrm{V}}$ wird in der Schaltung nach dem vorstehenden Bild umgesetzt?
        
        $P_\mathrm{V} = P_1 + P_2 + P_3 = U_\mathrm{R1} \cdot I_1 + U_\mathrm{R2} \cdot I_2 + U_\mathrm{R3} \cdot I_3$\

        $= U_\mathrm{Q} \cdot \frac{R_1}{R_1+R_2} \cdot \frac{U_\mathrm{Q}}{R_1+R_2} +U_\mathrm{Q} \cdot \frac{R_2}{R_1+R_2} \cdot \frac{U_\mathrm{Q}}{R_1+R_2} + R_3 \cdot I_\mathrm{Q} \cdot I_\mathrm{Q} $

        $= U_\mathrm{Q}^2 \cdot \frac{R_1}{(R_1+R_2)^2} + U_\mathrm{Q}^2 \cdot \frac{R_2}{(R_1+R_2)^2} + I_\mathrm{Q}^2 \cdot R_3 = \frac{U_\mathrm{Q}^2}{R_1+R_2} +I_\mathrm{Q}^2 \cdot R_3 $


        \item Die Klemmen A und B werden kurzgeschlossen. Welcher Strom $I_{\mathrm{K}}$ fließt durch diese Verbindung
        zwischen den Klemmen A und B?

                \begin{center}
            \begin{tikzpicture}
    \draw(3,3) to[short,o-](2,3)
    to[R, i<_,l_=$R_1$, name = R1](0,3)
    to[V,v_,name=UQ](0,0)
    to[short](7,0)
    to[V,v_,name=IQ](7,3)
    to[R,l=$R_3$,-o](4,3)
    to[short,i<, name=IK](3,3);

    \draw(2.5,0)to[R, i_<, l=$R_2$,*-*, name = R2](2.5,3);
    %\draw(5,0)to[R,l=$R_3$,*-*](5,3);
    \iarrmore{R1}{$I_{\mathrm{1}}$};
    \iarrmore{R2}{$I_{\mathrm{2}}$};
    \iarrmore{IK}{$I_{\mathrm{K}}$};
    \varrmore{UQ}{$U_{\mathrm{Q}}$};
    \varrmore{IQ}{$I_{\mathrm{Q}} \cdot R_3$};
    

    \node at (3.2,3.3){A};
    \node at (3.8,3.3){B};
    \node[text=blue] at (1.2,1.5){M I};
    \node[text=blue] at (5.2,1.5){M II};
    \node[text=red] at (2.4,3.3){K I};
    %\node at (5.2,1.3){M II};
\end{tikzpicture}
        \end{center}

        MI: $R_1 \cdot I_1$ + $R_2 \cdot I_2$             - $U_0$ = 0\\
        MII: $- R_2 \cdot I_2 $ + $R_3 \cdot I_\mathrm{K}$      - $I_\mathrm{Q} \cdot R_3$ = 0\\
        KI: $I_1$ - $I_2$ - $I_\mathrm{K}$ = 0\\

        3 Gleichungen, 3 Unbekannte $\rightarrow$ Lösung über beispielsweise über Gauss-Algorithmus möglich.\\

        Einfacher:\\
        ESB aus c) 

        \begin{center}
            \input{Tikz/src/Aufgaben/Kirchhoff_A4_lsg3}
        \end{center}

        $U_0 = U_\mathrm{AB,L} =  U_\mathrm{Q} \cdot \frac{R_2}{R_1+R_2} + I_\mathrm{Q} \cdot R_3$\\

        $ R_\mathrm{i} = R_\mathrm{i,AB}= R_3 + \frac{R_1 \cdot R_2}{R_1 + R_2}$\\

        $\rightarrow I_\mathrm{K} = \frac{U_0}{R_\mathrm{i}}= \frac{U_\mathrm{Q}\cdot R_2 + I_\mathrm{Q} \cdot R_3 \cdot (R_1+R_2)}{R_1 \cdot R_2 + R_3 \cdot (R_1+R_2)} $






        \item Welche Spannung $U_{\mathrm{AB}}$ stellt sich ein, wenn an die Klemmen A und B ein Widerstand $R$ angeschlossen wird? $(R_1 = R_2 = R_3 = R)$    
    

    Aus c), Widerstand $R$ eingefügt
    \begin{center}
        \begin{tikzpicture}
    \draw(3,3) to[short,o-](2,3)
    to[R,l_=$R_\mathrm{i}$](0,3)
    to[V,v_,name=UQ](0,0)
    to[short, -o](3,0)
    to[R, l_=$R$](3,3);


    \varrmore{UQ}{$U_{\mathrm{0}}$};
  

    \node at (3.0,3.3){A};
    \node at (3.3,0){B};
\end{tikzpicture}
    \end{center}
    $R_\mathrm{i} = \frac{R \cdot R}{R + R} + R= \frac{1}{2} \cdot R + R = \frac{3}{2} \cdot R$\\

    Aus d): $U_0 = U_\mathrm{Q} \cdot \frac{R}{R+R} + R \cdot I_\mathrm{Q} = \frac{1}{2} U_\mathrm{Q} + R \cdot I_\mathrm{Q}$

    Spannungsteiler:\\
    
    $U_\mathrm{AB} = U_0 \cdot \frac{R}{2,5 \cdot R}$\\

    $ = 0,4 \cdot U_0 = U_\mathrm{Q} \cdot 0,5 \cdot 0,4 + I_\mathrm{Q} \cdot R \cdot 0,4 = 0,2 \cdot U_\mathrm{Q} + 0,4 \cdot R \cdot I_\mathrm{Q}$




    \end{enumerate}

}